\section{Limitations}
\label{sec:limitations}

\paragraph{Simulator realism.} The model is flow-level: delay and loss are
analytic functions of utilization, traffic is exogenous and inelastic, the
control plane acts instantaneously, fast reroute is free, and telemetry is
exact. Learned routing policies trained in fluid models are known not to
transfer to packet-level dynamics~\citep{boltres2024packerl}; our conclusions
are about this model. Inelastic traffic in particular removes a natural
source of temporal coupling (congestion that suppresses future demand).

\paragraph{Formulation.} One trunk may move per interval network-wide and each
demand has four fixed candidate paths. These choices create the incremental,
stateful character of the problem; a controller that re-optimizes all trunks
at once faces the one-shot problem studied by Teal and DOTE instead. We did
not vary the number of candidates or the per-interval move limit, because the
environment definition is frozen.

\paragraph{Scale and data.} One hand-designed 18-router topology, 17 demands,
hand-written diurnal profiles and seven scripted test scenarios; no operator
traffic matrices or failure traces. No claim of topology generalization is
made.

\paragraph{Learning.} Five training roots for the primary comparison and
two to three for ablations; root-level intervals are wide and are reported.
A budget of 400k transitions (about 1{,}390 simulated days), with one longer
run per learner. PPO received a bounded one-factor sensitivity study, the
bandit none; neither used graph-structured networks, observation
normalization or learning-rate schedules. A discount below the true one can
help with limited data~\citep{jiang2015horizon,amit2020discount}, so ``the
myopic learner wins at this budget'' does not imply that no sequential
structure exists---which is why we measure the structure separately.

\paragraph{Objective.} All rankings use one hand-weighted operational reward;
we report its components and operational metrics alongside, but a different
operator objective could reorder close methods.

\paragraph{References.} The clairvoyant oracles see future traffic and only
evaluate ``one move, then no change''; they bound $H$-step-greedy reasoning, not
the optimal policy. MILP-track optimizes gross maximum utilization only.

\paragraph{Provenance.} The closed study's checkpoints and per-episode
holdout data were not archived; its learner results are compared at the
aggregate level it committed.
